\section{Related work}\label{sec:related}
\citet{kim2024parity} prove efficient teacher-forced parity learning with a Transformer and compare learning with and without chain of thought. Their parity setting, prescribed nonlinearity and attention parameterization differ from this paper's biased modular increments, local fixed-ReLU analysis and supplied-routing family. Our contribution is the matched finite-time plateau and persistent transition learning in the specified two-block softmax/ReLU parameterization; it does not subsume parity learning at $p=1/2$.

Makkuva and collaborators analyze Transformer loss landscapes and initialization-dependent dynamics on Markov data \citep{makkuva2025attention,makkuva2024local}. Their Markov next-token setting motivates an architecture-level optimization analysis. Here, the chain is the hidden modular state, and the comparison changes whether its intermediate states are displayed. We use Fourier conditioning rather than identify a separate attracting local minimum.

Gradient barriers based on weak target correlations \citep{shalevshwartz2017failures,abbe2025parities} explain why accessible features can reveal little about a high-order answer. The additional work here is controlling letter support, Taylor remainders and parameter scale in the Transformer, then converting a local loss floor into a trajectory bound. Automata constructions \citep{liu2023automata} show how shallow Transformers can represent sequential shortcuts; they do not imply convergence within the trained subspace. Scratchpad supervision \citep{nye2021scratchpads}, task decomposition \citep{wies2023subtask}, and locality of experience \citep{prystawski2023think} provide related perspectives on intermediate states. Expressivity and statistical advantages of chain of thought \citep{merrill2024expressive,wen2024sparse} concern complementary questions.

